% =============================================================================
% Appendix F: Fault Injection Specifications
% ORIUS/DC3S Battery-Safety Thesis
% =============================================================================
\chapter{Fault Injection Specifications}
\label{app:fault-specs}

This appendix specifies the six fault families used throughout the
CPSBench battery track (Chapter~\ref{ch:cpsbench}).  Each family
models a distinct telemetry-degradation mechanism observed in real BMS
and SCADA deployments.  The combined profile \texttt{drift\_combo}
superimposes all six families in a single scenario.

\section{Fault Families}

Table~\ref{tab:fault-families} defines each fault family, its
configurable parameters, the tested range, and the battery-domain
interpretation.

\begin{longtable}{@{}p{2.0cm}p{3.0cm}p{2.5cm}p{5.5cm}@{}}
\caption{Battery fault families for CPSBench telemetry degradation.}
\label{tab:fault-families}\\
\toprule
\textbf{Fault} & \textbf{Parameters} & \textbf{Range tested} & \textbf{Battery interpretation} \\
\midrule
\endfirsthead
\multicolumn{4}{c}{\emph{(continued)}}\\
\toprule
\textbf{Fault} & \textbf{Parameters} & \textbf{Range tested} & \textbf{Battery interpretation} \\
\midrule
\endhead
\bottomrule
\endfoot

\texttt{dropout}
  & drop\_rate $d$
  & $d \in \{0.05,\allowbreak 0.10,\allowbreak 0.15,\allowbreak 0.20,\allowbreak 0.25,\allowbreak 0.30\}$
  & Packet loss on the BMS$\to$EMS telemetry link.  Dropped readings are replaced by the last received value (zero-order hold).
\\[4pt]

\texttt{delay\_jitter}
  & base delay $\tau_b$ (s), jitter $\sigma_\tau$ (s)
  & $\tau_b \in \{1,\allowbreak 3,\allowbreak 5,\allowbreak 10\}$
  & Variable network latency between BMS and the dispatch engine.  Delayed packets arrive out of order and are consumed at their arrival time.
\\[4pt]

\texttt{spike}
  & magnitude $\mu_s$ (multiplier)
  & $\mu_s \in \{2,\allowbreak 3,\allowbreak 4,\allowbreak 5\}$
  & Transient sensor glitch producing a reading $\mu_s$ times the true value.  Models ADC saturation or electromagnetic interference on current/voltage sensors.
\\[4pt]

\texttt{stale}
  & hold duration $h$ (steps)
  & $h \in \{2,\allowbreak 4,\allowbreak 6,\allowbreak 8,\allowbreak 12\}$
  & Sensor freeze: the telemetry value is held constant for $h$ consecutive steps.  Models BMS firmware hang or communication bus lock-up.
\\[4pt]

\texttt{drift}
  & drift rate $\beta$ (per step)
  & $\beta \in \{0.01,\allowbreak 0.02,\allowbreak 0.03,\allowbreak 0.05\}$
  & Slow, linear bias accumulating as $\xi_t = \beta\,t$.  Models sensor calibration drift (e.g.\ current-sensor offset) or gradual SOH-induced SOC bias.
\\[4pt]

\texttt{reorder}
  & buffer size $B$ (steps)
  & $B \in \{2,\allowbreak 4,\allowbreak 6,\allowbreak 8\}$
  & Out-of-order packet delivery: readings within a sliding window of $B$ steps are randomly permuted.  Models multi-path routing or priority-queue inversion on industrial networks.
\\
\end{longtable}

\section{Combined Profile: \texttt{drift\_combo}}

The \texttt{drift\_combo} profile activates all six fault families
simultaneously at their mid-range parameter values with a configurable
dropout rate $d$.  This is the default stress profile used in the main
results (Chapter~\ref{ch:main-results}) and in the extended controller
comparison (Appendix~\ref{app:extended-results}).

\begin{table}[htbp]
\centering
\caption{\texttt{drift\_combo} default parameter values.}
\label{tab:drift-combo-defaults}
\small
\begin{tabular}{@{}llr@{}}
\toprule
\textbf{Fault} & \textbf{Parameter} & \textbf{Default value} \\
\midrule
dropout        & $d$          & 0.20 \\
delay\_jitter  & $\tau_b$     & 3\,s \\
spike          & $\mu_s$      & 3$\times$ \\
stale          & $h$          & 4 steps \\
drift          & $\beta$      & 0.02\,/step \\
reorder        & $B$          & 4 steps \\
\bottomrule
\end{tabular}
\end{table}

\section{Fault Injection Implementation}

Faults are injected by the CPSBench runner
(\texttt{src/orius/cpsbench\_iot/runner.py}) between the plant's
true-state output and the controller's observation input.  The
injection pipeline:

\begin{enumerate}
  \item \textbf{Plant emits} clean telemetry $\zt$.
  \item \textbf{Fault layer} transforms $\zt \to \ztt$ according to the
    active fault families and their parameters.
  \item \textbf{Controller receives} degraded telemetry $\ztt$ as its
    observation $\ot$.
  \item \textbf{DC3S OQE} computes $\wt$ from $\ztt$ (detecting the
    degradation).
\end{enumerate}

\noindent
Fault parameters are set in \texttt{configs/dc3s.yaml} under the
\texttt{fault\_profile} key.  Each evaluation run logs the active
profile and parameter snapshot in the certificate metadata for
reproducibility.

\section{Severity Mapping}

Table~\ref{tab:fault-severity} maps each fault family to its
primary impact on DC3S pipeline stages.

\begin{table}[htbp]
\centering
\caption{Fault-to-pipeline severity mapping.}
\label{tab:fault-severity}
\small
\begin{tabular}{@{}lcccc@{}}
\toprule
\textbf{Fault}
  & \textbf{OQE $w_t$}
  & \textbf{RUI $m_t$}
  & \textbf{Shield IR}
  & \textbf{Cert.\ validity} \\
\midrule
dropout        & $\downarrow\downarrow$   & $\uparrow\uparrow$   & $\uparrow$   & $\downarrow$   \\
delay\_jitter  & $\downarrow$             & $\uparrow$            & $\uparrow$   & $\downarrow$   \\
spike          & $\downarrow\downarrow$   & $\uparrow\uparrow$   & $\uparrow\uparrow$ & $\downarrow\downarrow$ \\
stale          & $\downarrow$             & $\uparrow$            & $\uparrow$   & $\downarrow$   \\
drift          & $\downarrow\downarrow$   & $\uparrow\uparrow$   & $\uparrow$   & $\downarrow\downarrow$ \\
reorder        & $\downarrow$             & $\uparrow$            & $\uparrow$   & ---                     \\
\bottomrule
\end{tabular}
\medskip

\footnotesize
$\downarrow$ = mild decrease; $\downarrow\downarrow$ = strong decrease;
$\uparrow$ = mild increase; $\uparrow\uparrow$ = strong increase;
--- = negligible effect.
\end{table}
